\nonstopmode{}
\documentclass[letterpaper]{book}
\usepackage[times,inconsolata,hyper]{Rd}
\usepackage{makeidx}
\usepackage[utf8]{inputenc} % @SET ENCODING@
% \usepackage{graphicx} % @USE GRAPHICX@
\makeindex{}
\begin{document}
\chapter*{}
\begin{center}
{\textbf{\huge Package `deepMatchR'}}
\par\bigskip{\large \today}
\end{center}
\inputencoding{utf8}
\ifthenelse{\boolean{Rd@use@hyper}}{\hypersetup{pdftitle = {deepMatchR: Deep Learning Models for HLA Matching}}}{}
\ifthenelse{\boolean{Rd@use@hyper}}{\hypersetup{pdfauthor = {Nick Borcherding}}}{}
\begin{description}
\raggedright{}
\item[Title]\AsIs{Deep Learning Models for HLA Matching}
\item[Version]\AsIs{0.99.0}
\item[Date]\AsIs{2024-07-24}
\item[Author]\AsIs{Nick Borcherding [aut, cre]}
\item[Maintainer]\AsIs{Nick Borcherding }\email{ncborcherding@gmail.com}\AsIs{}
\item[Description]\AsIs{A set of tools and models for analyzing HLA matching. This includes the identification and quantification of immunogenicity mismatches, the generation of in silico patients based on frequency table, and interaction with common tools in HLA. }
\item[License]\AsIs{MIT + file LICENSE}
\item[Encoding]\AsIs{UTF-8}
\item[LazyData]\AsIs{true}
\item[LazyDataCompression]\AsIs{xz}
\item[RoxygenNote]\AsIs{7.3.2}
\item[biocViews]\AsIs{Software, Classification, Annotation, Sequencing}
\item[Depends]\AsIs{R (>= 4.0)}
\item[Imports]\AsIs{dplyr, ggplot2, grDevices, readxl, stringr, tidyr, data.table,
patchwork, directlabels, keras3, pracma, purrr, magrittr,
tibble, treemapify, immReferent}
\item[Suggests]\AsIs{spelling, testthat}
\item[Remotes]\AsIs{BorchLab/immReferent}
\item[Language]\AsIs{en-US}
\item[URL]\AsIs{}\url{https://github.com/ncborcherding/deepMatchR/}\AsIs{}
\item[BugReports]\AsIs{}\url{https://github.com/ncborcherding/deepMatchR/issues}\AsIs{}
\end{description}
\Rdcontents{\R{} topics documented:}
\inputencoding{utf8}
\HeaderA{calculateAuc}{Calculate Antigen AUC Based on MFI}{calculateAuc}
\aliasA{cregAuc}{calculateAuc}{cregAuc}
\aliasA{epletAuc}{calculateAuc}{epletAuc}
\aliasA{serologyAuc}{calculateAuc}{serologyAuc}
%
\begin{Description}
This function provides a unified interface for analyzing Single Antigen Bead
(SAB) data against different antigenic features, such as eplets,
cross-reactive groups (CREGs), or serology. It computes the proportion of features
that are positive above a sequence of MFI cut-offs and integrates that to
obtain an area-under-the-curve (AUC).

Depending on user arguments, it can either generate a ggplot or return
the AUC results as a tibble.
\end{Description}
%
\begin{Usage}
\begin{verbatim}
calculateAuc(
  result_file,
  analysis_type,
  feature_filter = 3,
  percPos_filter = 0.8,
  group_by = NULL,
  label = TRUE,
  cut_min = 250,
  cut_max = 10000,
  cut_step = 250,
  plot_results = TRUE,
  palette = "spectral",
  evidence_level = c("A1", "A2"),
  eplet_filter = 3,
  top_eplets = 10,
  creg_filter = 3,
  serology_filter = NULL,
  ...
)

epletAuc(
  result_file,
  evidence_level = c("A1", "A2"),
  group_by = "eplet",
  label = TRUE,
  eplet_filter = 3,
  percPos_filter = 0.8,
  cut_min = 250,
  cut_max = 10000,
  cut_step = 250,
  plot_results = TRUE,
  top_eplets = 10,
  palette = "spectral"
)

cregAuc(
  result_file,
  group_by = "creg",
  label = TRUE,
  creg_filter = 3,
  percPos_filter = 0.8,
  cut_min = 250,
  cut_max = 10000,
  cut_step = 250,
  plot_results = TRUE,
  palette = "spectral"
)

serologyAuc(result_file, serology_filter = 3, ...)
\end{verbatim}
\end{Usage}
%
\begin{Arguments}
\begin{ldescription}
\item[\code{result\_file}] A data frame of SAB results or a path to a CSV/XLS/XLSX file.

\item[\code{analysis\_type}] Character. The type of analysis to perform.
Must be one of `"eplet"`, `"creg"`, `"serology"`.

\item[\code{feature\_filter}] Integer. The minimum number of beads/alleles that must
carry the feature for it to be included in the analysis. Defaults to `3`.

\item[\code{percPos\_filter}] Numeric (0-1). The minimum proportion of beads that must
be positive for at least one cut-off to keep the feature. Defaults to `0.8`.

\item[\code{group\_by}] Character. The aesthetic used to color curves in the plot.

\item[\code{label}] Logical. If `TRUE`, labels the curves on the plot. Defaults to `TRUE`.

\item[\code{cut\_min, cut\_max, cut\_step}] Range and step for MFI cut-offs.

\item[\code{plot\_results}] Logical. If `TRUE` (default), returns a `ggplot` object.
If `FALSE`, returns a summarized tibble with AUC results.

\item[\code{palette}] Character. A color palette name (see `grDevices::hcl.pals`)
or a custom palette function. Defaults to `"spectral"`.

\item[\code{evidence\_level}] For eplet analysis, a character vector of evidence levels to keep.

\item[\code{eplet\_filter}] For eplet analysis, the minimum number of times an eplet must appear.

\item[\code{top\_eplets}] For eplet analysis, the maximum number of top eplets to display.

\item[\code{creg\_filter}] For CREG analysis, the minimum number of times a CREG must appear.

\item[\code{serology\_filter}] For serology analysis, a filter to be applied.

\item[\code{...}] Additional arguments passed to the plot theme.
\end{ldescription}
\end{Arguments}
%
\begin{Value}
Either a `ggplot` object or a tibble with AUC results. The tibble
will contain columns for the feature (`eplet`, `creg`, `serology`), `AUC`,
`norm\_AUC`, `total\_count`, and `loci`.
\end{Value}
\inputencoding{utf8}
\HeaderA{compareHlaSequences}{Compare two HLA protein sequences and identify polymorphisms}{compareHlaSequences}
%
\begin{Description}
This function compares two HLA protein sequences of the same length and
identifies the positions at which the amino acids differ. It also
characterizes the change in physicochemical properties (polarity and charge)
at each polymorphic site based on IMGT classifications.
\end{Description}
%
\begin{Usage}
\begin{verbatim}
compareHlaSequences(seq1, seq2)
\end{verbatim}
\end{Usage}
%
\begin{Arguments}
\begin{ldescription}
\item[\code{seq1}] A character string representing the first protein sequence.

\item[\code{seq2}] A character string representing the second protein sequence.
\end{ldescription}
\end{Arguments}
%
\begin{Value}
A data frame with the following columns for each polymorphism:
\begin{itemize}

\item{} \code{Position}: The position of the polymorphism (1-based index).
\item{} \code{AA\_Seq1}: The amino acid in sequence 1 at the polymorphic site.
\item{} \code{AA\_Seq2}: The amino acid in sequence 2 at the polymorphic site.
\item{} \code{Polarity\_Change}: A string describing the change in polarity (e.g., "Polar to Nonpolar").
\item{} \code{Charge\_Change}: A string describing the change in charge (e.g., "Positive to Negative").

\end{itemize}

If the sequences are identical, it returns an empty data frame.
\end{Value}
%
\begin{Examples}
\begin{ExampleCode}
seq1 <- "YFAMYGEKVAHTHVDTLYVRYHY"
seq2 <- "YFDMYGEKVAHTHVDTLYVRYHY"
compareHlaSequences(seq1, seq2)
\end{ExampleCode}
\end{Examples}
\inputencoding{utf8}
\HeaderA{deepMatchR\_cregs}{CREG-Allele Mapping Data (Dummy)}{deepMatchR.Rul.cregs}
\keyword{datasets}{deepMatchR\_cregs}
%
\begin{Description}
A dummy dataset containing a mapping of HLA alleles to their
Cross-Reactive Groups (CREGs). This is a placeholder file to allow the
package to build and run correctly. The data is not comprehensive or
scientifically validated.
\end{Description}
%
\begin{Usage}
\begin{verbatim}
deepMatchR_cregs
\end{verbatim}
\end{Usage}
%
\begin{Format}
A data frame with 6 rows and 2 variables:
\begin{description}

\item[allele] An HLA allele string.
\item[creg] The CREG group corresponding to the allele.

\end{description}

\end{Format}
\inputencoding{utf8}
\HeaderA{deepMatchR\_eplets}{Eplet information extracted from [HLA Eplet Registry](https://www.epregistry.com.br/). Eplets were filtered to retain evidence classes A1, A2, B and D).}{deepMatchR.Rul.eplets}
%
\begin{Description}
Eplet information extracted from [HLA Eplet Registry](https://www.epregistry.com.br/).
Eplets were filtered to retain evidence classes A1, A2, B and D).
\end{Description}
\inputencoding{utf8}
\HeaderA{deepMatchR\_example}{Simulated single antigen bead (SAB) assay for class I and class II and PRA assay with fixed formatting.}{deepMatchR.Rul.example}
%
\begin{Description}
Simulated single antigen bead (SAB) assay
for class I and class II and PRA assay
with fixed formatting.
\end{Description}
\inputencoding{utf8}
\HeaderA{getAlleleSequence}{Get Amino Acid Sequence for an HLA Allele}{getAlleleSequence}
%
\begin{Description}
This function retrieves the amino acid sequence for a given HLA allele
name using the `immReferent` package.
\end{Description}
%
\begin{Usage}
\begin{verbatim}
getAlleleSequence(allele_name)
\end{verbatim}
\end{Usage}
%
\begin{Arguments}
\begin{ldescription}
\item[\code{allele\_name}] A character string representing the HLA allele name
(e.g., "A*01:01").
\end{ldescription}
\end{Arguments}
%
\begin{Value}
A character string representing the amino acid sequence.
\end{Value}
\inputencoding{utf8}
\HeaderA{hlaGeno}{Create an hla\_genotype object}{hlaGeno}
\keyword{internal}{hlaGeno}
%
\begin{Description}
Creates a new `hla\_genotype` object from a data frame of HLA calls. The
object is a list containing the genotype data and a record of the loci
present.
\end{Description}
%
\begin{Usage}
\begin{verbatim}
hlaGeno(df)
\end{verbatim}
\end{Usage}
%
\begin{Arguments}
\begin{ldescription}
\item[\code{df}] A data frame where rows are individuals and columns represent
HLA alleles (e.g., A\_1, A\_2, B\_1, B\_2, ...).
\end{ldescription}
\end{Arguments}
%
\begin{Value}
An object of class `hla\_genotype`.
\end{Value}
\inputencoding{utf8}
\HeaderA{params}{SAB Model Parameters}{params}
%
\begin{Description}
A data object containing the transformation recipe for Class I
and Class II MFI values. This includes the method and for each allele, the
ceiling, mu (mean), and standard deviation values.
\end{Description}
%
\begin{Format}
A list with two elements, one for each class.
\end{Format}
\inputencoding{utf8}
\HeaderA{plotAntibodies}{Plot Antibody Data with Optional Antigen-Level Table or Time-Series Trend}{plotAntibodies}
\aliasA{plotPRA}{plotAntibodies}{plotPRA}
\aliasA{plotSAB}{plotAntibodies}{plotSAB}
%
\begin{Description}
This function generates a bar plot of SAB (Single Antigen Beads) or PRA (Panel-Reactive Antibody) results
from a provided data frame or a file path. It can also plot MFI values over time.
\end{Description}
%
\begin{Usage}
\begin{verbatim}
plotAntibodies(
  result_file,
  type = "SAB",
  class = "I",
  plot_trend = FALSE,
  bead_cutoffs = NULL,
  highlight_threshold = 2000,
  vline_dates = NULL,
  add_table = TRUE,
  x_text_angle = 90,
  palette = "spectral",
  highlight_antigen = NULL,
  ...
)

plotSAB(..., type = "SAB")

plotPRA(..., type = "PRA")
\end{verbatim}
\end{Usage}
%
\begin{Arguments}
\begin{ldescription}
\item[\code{result\_file}] A data frame, a list of data frames (for trend plot), or a character string
specifying the path to a file in CSV, XLS, or XLSX format.

\item[\code{type}] Character. The type of assay, either "SAB" or "PRA". Defaults to "SAB".

\item[\code{class}] Character. For PRA plots, the class of the assay, either "I" or "II". Defaults to "I".

\item[\code{plot\_trend}] Logical. If TRUE, a time-series plot is generated. Defaults to FALSE.

\item[\code{bead\_cutoffs}] Numeric vector. Cutoff values for categorizing MFI values.
Defaults to `c(2000, 1000, 500, 250)` for SAB and `c(1500, 1000, 500, 250)` for PRA.

\item[\code{highlight\_threshold}] Numeric. MFI threshold for highlighting alleles in the trend plot. Defaults to 2000.

\item[\code{vline\_dates}] Vector of dates. Dates to draw vertical lines on the trend plot.

\item[\code{add\_table}] Logical. Whether to add the antigen-level information as a
table to the bottom of the bar plot. Defaults to TRUE.

\item[\code{x\_text\_angle}] Numeric. Angle for the antigen/allele text in the table. Defaults to 90.

\item[\code{palette}] Character. A color palette name. Defaults to "spectral".

\item[\code{highlight\_antigen}] Character vector. Optional antigen(s) to highlight.

\item[\code{...}] Additional arguments passed to the ggplot theme.
\end{ldescription}
\end{Arguments}
%
\begin{Value}
A `ggplot` object.
\end{Value}
\inputencoding{utf8}
\HeaderA{plotEplets}{Plot Eplet Results from SPI Assay}{plotEplets}
%
\begin{Description}
This function reads in SAB (single antigen bead) results (either as a data frame or
as a file path to a CSV/XLS/XLSX file), processes the data to quantify eplet-specific
positivity based on a specified MFI cutoff, and then generates one of three plot types:
a treemap, a bar plot, or an AUC plot. The eplet annotations are joined from an internal
database, and the resulting plot is colored by evidence level.
\end{Description}
%
\begin{Usage}
\begin{verbatim}
plotEplets(
  result_file,
  cutoff = 2000,
  evidence_level = c("A1", "A2"),
  group_by = "eplet",
  eplet_filter = 3,
  percPos_filter = 0.4,
  plot_type = c("treemap", "bar", "AUC"),
  cut_min = 250,
  cut_max = 10000,
  cut_step = 250,
  top_eplets = 10,
  palette = "spectral",
  ...
)
\end{verbatim}
\end{Usage}
%
\begin{Arguments}
\begin{ldescription}
\item[\code{result\_file}] A data frame containing SAB results or a character string specifying
the path to a SAB file in CSV, XLS, or XLSX format.

\item[\code{cutoff}] Numeric. Threshold for MFI to separate positive and negative beads.
Default is 2000.

\item[\code{evidence\_level}] Character vector indicating the antibody reactivity levels to keep.
Defaults to \code{c("A1", "A2")}, which represent antibody-confirmed eplets.
Other acceptable levels include \code{"B"}, \code{"D"}, or \code{NULL} to apply no filter.

\item[\code{group\_by}] A character string or indicating the coloring grouping for
the plot, default is `eplet`. Other options include `loci` or
`evidence\_level`.

\item[\code{eplet\_filter}] Integer. (Only used when \code{plot\_type = "AUC"})
Specifying the minimum number of times an eplet must appear in the assay
before calculating the AUC. Defaults to `3`.

\item[\code{percPos\_filter}] Numeric value. (Only used when \code{plot\_type = "AUC"})
Value between 0 and 1 specifying the minimum relative proportion of
positive beads (per eplet) to include in the final summary. Eplets with a
proportion below this threshold are excluded. Defaults to 0.4.

\item[\code{plot\_type}] Character. Type of plot to generate. Must be one of \code{"treemap"},
\code{"bar"}, or \code{"AUC"}. Defaults to \code{"treemap"}.

\item[\code{cut\_min}] Integer. (Only used when \code{plot\_type = "AUC"}) The minimum MFI cutoff value.
Defaults to 250.

\item[\code{cut\_max}] Integer. (Only used when \code{plot\_type = "AUC"}) The maximum MFI cutoff value.
Defaults to 10000.

\item[\code{cut\_step}] Integer. (Only used when \code{plot\_type = "AUC"}) The increment step between
the minimum and maximum MFI cutoff values. Defaults to 250.

\item[\code{top\_eplets}] Integer. The maximum number of top eplets to display in the bar or AUC plot.
Defaults to 10.

\item[\code{palette}] Character. A color palette name (from \LinkA{hcl.pals}{hcl.pals}) or a custom
palette function to use for the plot. Defaults to \code{"spectral"}.

\item[\code{...}] Additional arguments passed to the ggplot theme
\end{ldescription}
\end{Arguments}
%
\begin{Value}
A \code{ggplot} object visualizing eplet counts (or AUC values) according to the
specified parameters.
\end{Value}
%
\begin{Examples}
\begin{ExampleCode}
# Using a data frame:
plotEplets(deepMatchR_example[[1]],
           cutoff = 2000,
           evidence_level = c("A1", "A2", "B"),
           percPos_filter = 0.4,
           plot_type = "treemap")
\end{ExampleCode}
\end{Examples}
\inputencoding{utf8}
\HeaderA{print.hla\_genotype}{Print an hla\_genotype object}{print.hla.Rul.genotype}
%
\begin{Description}
Print an hla\_genotype object
\end{Description}
%
\begin{Usage}
\begin{verbatim}
## S3 method for class 'hla_genotype'
print(x, ...)
\end{verbatim}
\end{Usage}
%
\begin{Arguments}
\begin{ldescription}
\item[\code{x}] An object of class `hla\_genotype`.

\item[\code{...}] Additional arguments (not used).
\end{ldescription}
\end{Arguments}
%
\begin{Value}
Invisibly returns the original object.
\end{Value}
\inputencoding{utf8}
\HeaderA{quantifyEpletMismatch}{Quantify Eplet Mismatches Between Two Alleles}{quantifyEpletMismatch}
%
\begin{Description}
This function calculates the number of eplet mismatches between two HLA
alleles. It uses the internal `deepMatchR\_eplets` dataset.
\end{Description}
%
\begin{Usage}
\begin{verbatim}
quantifyEpletMismatch(allele1, allele2)
\end{verbatim}
\end{Usage}
%
\begin{Arguments}
\begin{ldescription}
\item[\code{allele1}] A character string for the first HLA allele (e.g., "A*01:01").

\item[\code{allele2}] A character string for the second HLA allele (e.g., "A*02:01").
\end{ldescription}
\end{Arguments}
%
\begin{Value}
An integer representing the number of mismatched eplets.
\end{Value}
\inputencoding{utf8}
\HeaderA{quantifyMismatch}{Quantify Amino Acid Mismatches Between Two Sequences}{quantifyMismatch}
%
\begin{Description}
This function compares two amino acid sequences of the same length and
returns the number of positions at which the amino acids differ.
\end{Description}
%
\begin{Usage}
\begin{verbatim}
quantifyMismatch(sequence1, sequence2)
\end{verbatim}
\end{Usage}
%
\begin{Arguments}
\begin{ldescription}
\item[\code{sequence1}] A character string representing the first amino acid sequence.

\item[\code{sequence2}] A character string representing the second amino acid sequence.
\end{ldescription}
\end{Arguments}
%
\begin{Value}
An integer representing the number of mismatches.
\end{Value}
%
\begin{Examples}
\begin{ExampleCode}
seq1 <- "YFAMYGEKVAHTHVDTLYVRYHY"
seq2 <- "YFDMYGEKVAHTHVDTLYVRFHY"
quantifyMismatch(seq1, seq2)
#> [1] 2

\end{ExampleCode}
\end{Examples}
\inputencoding{utf8}
\HeaderA{validateHlaGeno}{Validate an hla\_genotype object}{validateHlaGeno}
\keyword{internal}{validateHlaGeno}
%
\begin{Description}
Validate an hla\_genotype object
\end{Description}
%
\begin{Usage}
\begin{verbatim}
validateHlaGeno(x)
\end{verbatim}
\end{Usage}
%
\begin{Arguments}
\begin{ldescription}
\item[\code{x}] An object to validate.
\end{ldescription}
\end{Arguments}
%
\begin{Value}
`TRUE` if the object is a valid `hla\_genotype` object, otherwise
throws an error.
\end{Value}
\printindex{}
\end{document}
